\documentclass[12pt]{article}
\usepackage[margin=1in]{geometry}
\usepackage[utf8]{inputenc}
\usepackage{latexsym, amsfonts, enumitem, graphicx, environ,enumitem, fancyhdr, color, amssymb, amsthm, amsmath, mathtools, tikz, nicematrix}
\usetikzlibrary{patterns}
\usetikzlibrary{decorations.pathreplacing}
\pagestyle{fancy}
\setlength{\headheight}{65pt}
\newenvironment{problem}[2][Problem]{\begin{trivlist}
    \item[\hskip \labelsep {\bfseries #1}\hskip \labelsep {\bfseries #2.}]}{\end{trivlist}}
\DeclareMathOperator{\Ima}{Im}

\usepackage{hyperref}
\usepackage[capitalize]{cleveref}

\newtheorem{thm}{Theorem}[section]
\newtheorem{lem}[thm]{Lemma}
\newtheorem{cl}[thm]{Claim}
\newtheorem{prop}[thm]{Proposition}
\newtheorem{cor}[thm]{Corollary}
\newtheorem{conj}[thm]{Conjecture}
\newtheorem{obstruction}[thm]{Obstruction}

\usepackage{mdframed}

\theoremstyle{definition}
\newtheorem{definition}[thm]{Definition}
\newtheorem{remark}[thm]{Remark}
\newtheorem{question}{Question}
\newtheorem{example}[thm]{Example}
\newtheorem{exercise}[thm]{Exercise}
\newtheorem{properties}[thm]{Properties}

\usepackage{tcolorbox}
\tcbuselibrary{theorems,skins,breakable}

\tcolorboxenvironment{example}{
enhanced jigsaw,colframe=cyan,interior hidden,
breakable,%before skip=10pt,after skip=10pt
}

\tcolorboxenvironment{thm}{
enhanced jigsaw,colframe=red,interior hidden,
breakable,%before skip=10pt,after skip=10pt
}

\tcolorboxenvironment{prop}{
enhanced jigsaw,colframe=red,interior hidden,
breakable,%before skip=10pt,after skip=10pt
}

\tcolorboxenvironment{properties}{
enhanced jigsaw,colframe=red,interior hidden,
breakable,%before skip=10pt,after skip=10pt
}

\tcolorboxenvironment{lem}{
enhanced jigsaw,colframe=red,interior hidden,
breakable,%before skip=10pt,after skip=10pt
}
\tcolorboxenvironment{cor}{
enhanced jigsaw,colframe=red,interior hidden,
breakable,%before skip=10pt,after skip=10pt
}
\tcolorboxenvironment{obstruction}{
enhanced jigsaw,colframe=red,interior hidden,
breakable,%before skip=10pt,after skip=10pt
}

\tcolorboxenvironment{proof}{
enhanced jigsaw, frame hidden,%interior hidden,
breakable,%before skip=10pt,after skip=10pt
}

\renewcommand\qedsymbol{\( \blacksquare \)}

\lhead{Avi Chetlin \\ Assignment 02 \\ 09 September 2026}
\rhead{Dr. David Singer \\ MATH 465 Differential Geometry \\ Fall 2026}

\begin{document}
Unless explicitly stated, \(\alpha: I \rightarrow \mathbb{R}^3\) is a curve parametrized by arc length, \( s \), with curvature \( k(s) \neq 0 \), for all \( s \in I \).

\begin{problem}{1.5.1}
  Given the parametrized curve (helix)
  \begin{equation}
  \label{eq:1}
  \alpha(s) = \left( a \cos \frac{s}{c}, a \sin \frac{s}{c}, b \frac{s}{c} \right), \quad s \in \mathbb{R},
  \end{equation}
  where \( c^2 = a^2 + b^2 \),

  \begin{enumerate}[label=(\alph*)]
    \item Show that the parameter \( s \) is the arc length.
          \begin{proof}
            To prove this, we show that for arbitrary \( s \in I \), the arc length satisfies
            \begin{equation}
            \label{eq:3}
            \int\limits_0^s \left| \alpha^{\prime}(t) \right| \,dt = s.
            \end{equation}
            First of all,
            \begin{equation}
            \label{eq:4}
            \alpha^{\prime}(s) = \left( -\frac{a}{c}\sin \frac{s}{c}, \frac{a}{c}\cos \frac{s}{c}, \frac{b}{c} \right),
            \end{equation}
            so
            \begin{equation}
            \label{eq:5}
            \left| \alpha^{\prime} \right|^2 = \left( \frac{a}{c} \right)^2 + \left( \frac{b}{c} \right)^2 = \frac{a^2 + b^2}{c^2} = 1.
            \end{equation}
            Thus, our arc length integral becomes
            \begin{equation}
            \label{eq:6}
            \int\limits_0^s dt \equiv s.
            \end{equation}
          \end{proof}
    \item Determine the curvature and the torsion of \(\alpha\).
          \begin{proof}
            Since \(\alpha\) is parametrized by arc length, we have that \( k(s) = \left| \alpha^{\prime\prime}(s) \right| \).

            From \cref{eq:4},
            \begin{equation}
            \label{eq:7}
            \alpha^{\prime\prime}(s) = \left( - \frac{a}{c^2}\cos \frac{s}{c}, - \frac{a}{c^2}\sin \frac{s}{c}, 0 \right),
            \end{equation}
            and hence
            \begin{equation}
            \label{eq:8}
            \left| \alpha^{\prime\prime}(s) \right|^2 = \frac{a^2}{c^{4}},
            \end{equation}
            giving
            \begin{equation}
            \label{eq:9}
            \boxed{k(s) = \frac{a}{c^2}}.
            \end{equation}

            For the torsion, \( \tau \), recall the unit normal is defined by \( \alpha^{\prime\prime}(s) = k(s)N(s) \), so
            \begin{equation}
            \label{eq:10}
            N(s) = \frac{\alpha^{\prime\prime}(s)}{k(s)} = \left( -\cos \frac{s}{c}, -\sin \frac{s}{c}, 0 \right).
            \end{equation}

            Now, the Frenet equations tell us that \( B^{\prime} = \tau N \), where \( B = T \times N \); by application of the Leibnitz rule it is clear that \( B^{\prime} = T \times N^{\prime} \).
            We calculate
            \begin{align}
            \label{eq:12}
              B^{\prime} &= T \times N^{\prime} \\
              & = \left( \frac{b}{c^2}\cos \frac{s}{c}, \frac{b}{c^2} \sin \frac{s}{c}, 0 \right).
            \end{align}
            Moreover, \( B^{\prime} \cdot N = \tau \), hence
            \begin{equation}
            \label{eq:13}
            \boxed{\tau = - \frac{b}{c^2}}.
            \end{equation}
          \end{proof}
    \item Determine the osculating plane of \(\alpha\).
          \begin{proof}
            The osculating plane is spanned by \( T \) and \( N \).
            Let \( \mathbf{x} \in \mathbb{R}^3\) be an arbitrary point on the osculating plane of \(\alpha\).
            Since the binormal vector is, by definition, orthogonal to the unit tangent and normal vectors, we can write
            \begin{align}
            \label{eq:14}
              0 &= (\mathbf{x} - \alpha(s)) \cdot B(s) \\
                &= (\mathbf{x}^1 - a\cos \frac{s}{c}, \mathbf{x}^2 - a \sin \frac{s}{c}, \mathbf{x}^3 - b \frac{s}{c}) \cdot (\frac{b}{c}\sin \frac{s}{c}, - \frac{b}{c}\cos \frac{s}{c}, \frac{a}{c}) \\
                &= \frac{b}{c}\sin \frac{s}{c}\left( \mathbf{x}^1 - a \cos \frac{s}{c} \right) - \frac{b}{c}\cos \frac{s}{c}\left( \mathbf{x}^2 - a\sin \frac{s}{c} \right) + \frac{a}{c} \left( \mathbf{x}^3 - b \frac{s}{c} \right) \\
                &= \frac{b}{c}\left(\mathbf{x}^1 \sin \frac{s}{c} - \mathbf{x}^2 \cos \frac{s}{c}\right) + \frac{a}{c}\left( \mathbf{x}^3 - b \frac{s}{c} \right).
            \end{align}
            Multiplying both sides by \( c \) and rearranging, we obtain the equation for the osculating plane at the point \(\alpha(s)\),
            \begin{equation}
            \label{eq:15}
            \mathbf{x}^1b\sin \frac{s}{c} - \mathbf{x}^2b \cos \frac{s}{c} + \mathbf{x}^3 a = \frac{abs}{c}.
            \end{equation}
          \end{proof}
    \item Show that the lines containing \( N(s) \) and passing through \(\alpha(s)\) meet the \( z \)-axis under a constant angle equal to \( \frac{\pi}{2} \).
          \begin{proof}
            Since \( N(s) \) is a unit vector, we know that \( \left| \langle N(s), \hat{z} \rangle \right| = |N||\hat{z}|\cos\theta = \cos\theta \), where \(\theta\) is the angle between \( N \) and \( \hat{z}\).
            Since every line containing \( N(s) \) must meet the \( z \)-axis at the same angle, it suffices to show that \( \cos\theta \equiv 0 \).
            From \cref{eq:10},
            \begin{equation}
            \label{eq:16}
            \langle N(s), (0,0,1) \rangle = 0 = \cos\theta,
            \end{equation}
            hence \(\theta \equiv \frac{\pi}{2}\).

            Finally, these lines have the form
            \begin{equation}
            \label{eq:51}
            \alpha(s) + tN(s) = (\cos \frac{s}{c}(a-t), \sin \frac{s}{c}(a-t), b \frac{s}{c}) = (0,0, b \frac{s}{c}),
            \end{equation}
            which clearly intersects the \( z \)-axis when \( t = a \).
          \end{proof}
    \item Show that the tangent lines to \(\alpha\) make a constant angle with the \( z \)-axis.
          \begin{proof}
            It again suffices to show that \( \left| \langle T, (0,0,1) \rangle \right| = \text{const} \).
              With \cref{eq:4},
            \begin{equation}
            \label{eq:17}
            \cos\theta = \langle T, (0,0,1) \rangle = \frac{b}{c},
            \end{equation}
            hence \( T \) makes a constant angle \(\theta = \arccos \frac{b}{c}\) with \( \hat{z} \).
          \end{proof}
  \end{enumerate}
\end{problem}

\begin{problem}{1.5.14}
  Let \(\alpha:(a,b) \rightarrow \mathbb{R}^2\) be a regular parametrized plane curve.
  Assume that there exists \( t_0 \), \( a < t_0 < b \), such that the distance \( \left| \alpha(t) \right| \) from the origin to the trace of \(\alpha\) will be maximum at \( t_{0} \).
  Prove that the curvature \( k \) of \(\alpha\) at \( t_0 \) satisfies \( \left| k(t_0) \right| \geq 1/\left| \alpha(t_0) \right|\).
\end{problem}
\begin{proof}
  Since re-parametrizing \( \alpha \) preserves the curvature and norm of \(\alpha\), we may freely assume \(\alpha\) is parametrized by arc length.

  \(\left| \alpha(t) \right|\) taking a maximum at \( t=t_0 \) is equivalent to \( \left| \alpha(t) \right|^2 = \langle \alpha, \alpha \rangle \) taking a maximum, we must have
  \begin{equation}
  \label{eq:20}
  \frac{\mathrm{d}^2 }{\mathrm{d} t^2} \langle \alpha(t_0), \alpha(t_{0}) \rangle \leq 0.
  \end{equation}
  Hence, evaluating each term at \( t=t_0 \) but supressing arguments for the time being,
  \begin{align}
  \label{eq:21}
    0 &\geq 2\frac{\mathrm{d} }{\mathrm{d} t} \left( \langle\alpha^{\prime}, \alpha \rangle \right) \\
      &= 2 \left( \langle\alpha^{\prime\prime}, \alpha \rangle + \langle \alpha^{\prime}, \alpha^{\prime} \rangle \right) \\
      &= 2 \left( \langle T^{\prime}, \alpha \rangle + | T |^2 \right) = 2\langle kN , \alpha \rangle + 2 \\
    &= 2k \langle N , \alpha \rangle + 2,
  \end{align}
  and thus \(-1 \geq k \langle N, \alpha \rangle\), or
  \begin{equation}
  \label{eq:25}
  1 \leq k \langle -N, \alpha \rangle \leq |k| |N| |\alpha|
  \end{equation}
  which, recalling that \( |N| = 1 \), shows that
  \begin{equation}
  \label{eq:26}
  \frac{1}{|\alpha(t_0)|} \leq |k(t_0)|.
  \end{equation}
\end{proof}

\newpage
\begin{problem}{1.5.17}
  In general, a curve \(\alpha\) is called a \textit{helix} if the tangent lines of \(\alpha\) make a constant angle with a fixed direction.
  Assume that \(\tau(s) \neq 0\), \( s \in I \), and prove that
  \begin{enumerate}[label=(\alph*)]
    \item \(\alpha\) is a helix if and only if \( \frac{k}{\tau} \) is constant.
          \begin{proof}
            (\(\Rightarrow\)) Suppose \( \alpha \) is a helix, and let \( \hat{a} \) be the unit vector in the fixed direction with which the tangent lines make a constant angle \(\theta\).
            Then \( \langle T, \hat{a} \rangle = \cos\theta \) is constant, so
            \begin{equation}
            \label{eq:27}
            \langle T^{\prime}, \hat{a} \rangle + \langle T, \hat{a}^{\prime} \rangle = \langle T^{\prime}, \hat{a} \rangle = 0.
            \end{equation}
            Using the Frenet equations,
            \( \langle kN, \hat{a} \rangle = 0 \), so we can express \( \hat{a} \) in the Frenet frame as
            \begin{equation}
            \label{eq:28}
            \hat{a} = \cos\theta T + \beta B,
            \end{equation}
            where \( \beta \) is some constant; however, since \( \hat{a} \) is a unit vector, we must in particular have
            \begin{equation}
            \label{eq:29}
            \hat{a} = \cos\theta T + \sin\theta B.
            \end{equation}

            Differentiating \cref{eq:29} gives
            \begin{equation}
            \label{eq:30}
            0 = \cos\theta (kN) + \sin\theta (\tau N),
            \end{equation}
            which implies \(k = -\tau\tan\theta\), or
            \begin{equation}
            \label{eq:32}
            \frac{k}{\tau} = -\tan\theta = \text{constant}.
            \end{equation}

            (\(\Leftarrow\)) Suppose \( \frac{k}{\tau} \) is constant.
            Then \(\exists \theta\) such that \( -\tan\theta = \frac{k}{\tau}\), namely, \(\theta = \arctan -\frac{k}{\tau}\).
            This then implies that
            \begin{equation}
            \label{eq:35}
            k \cos\theta + \tau\sin\theta = 0.
            \end{equation}
            Multiplying both sides by \( N \) and using the Frenet equations gives
            \begin{equation}
            \label{eq:36}
            T^{\prime} \cos\theta + B^{\prime} \sin\theta = \frac{\mathrm{d} }{\mathrm{d} t} (T \cos\theta + B \sin\theta) = 0,
            \end{equation}
            so \( T\cos\theta + B \sin\theta \eqcolon \hat{a} \) is a constant unit vector.
            Moreover, \( \langle T, \hat{a} \rangle = \cos\theta \), which is constant since \(\theta\) is.
          \end{proof}
    \item \(\alpha\) is a helix if and only if the lines containing \( n(s) \) and passing through \( \alpha(s) \) are parallel to a fixed plane.
          \begin{proof}
            (\(\Rightarrow\)) Suppose \( \alpha \) is a helix, and let \( \hat{a} \) be the fixed unit vector such that \( \langle T, \hat{a} \rangle = \cos\theta \) is constant.
            We want to show there exists some plane \( \mathbf{P} \) with unit normal vector \( \hat{p} \) such that \( N, \hat{p} \rangle = 0 \).

            Notice that by being a helix, \(\alpha\) satisfies \( \langle T^{\prime}, \hat{a} \rangle = 0 \), that is, the unit normal of \(\alpha\) is orthogonal to \( \hat{a} \).
            Hence, \( N \) is parallel to a plane with normal vector \( \hat{a} \).

            (\(\Leftarrow\)) Suppose \( N \) is parallel to some fixed plane \( \mathbf{P} \).
            Then, letting \( \hat{a} \) be the vector normal to \( \mathbf{P} \), we have \( \langle N, \hat{a} \rangle = 0 \).
            Since
            \begin{equation}
            \label{eq:39}
            \langle N, \hat{a} \rangle = \frac{\mathrm{d} }{\mathrm{d} t} \langle T, \hat{a} \rangle = 0,
            \end{equation}
            we see that \( \langle T, \hat{a} \rangle = \text{constant} \), showing that \(\alpha\) is a helix.
          \end{proof}
    \item \(\alpha\) is a helix if and only if the lines containing \( b(s) \) and passing through \(\alpha(s)\) make a constant angle with a fixed direction.
          \begin{proof}
            (\(\Rightarrow\)) Suppose \(\alpha\) is a helix and let \( \hat{a} \) be the unit vector in the fixed direction with which \( T \) makes a constant angle \(\theta\).
            From part (a), we know we can write \( \hat{a} = \cos\theta T + \sin\theta B \), and therefore \( \langle B, \hat{a} \rangle = \sin \theta = \text{constant} \).

            (\(\Leftarrow\)) Suppose \( \langle B, \hat{a} \rangle = \text{constant} \) for some fixed direction \( \hat{a} \).
            Then, differentiating and applying the Frenet equations gives
            \begin{equation}
            \label{eq:42}
            \langle B^{\prime}, \hat{a} \rangle = \langle \tau N, \hat{a} \rangle = \tau \langle N , \hat{a} \rangle = 0,
            \end{equation}
            and since \( \frac{\mathrm{d} }{\mathrm{d} t} \langle T, \hat{a} \rangle = k \langle N, \hat{a} \rangle = 0 \), and \( k \neq 0 \), we must have that
            \begin{equation}
            \label{eq:43}
            \langle T , \hat{a} \rangle = \text{constant},
            \end{equation}
            hence \( \alpha \) is a helix.
          \end{proof}
    \item The curve
          \begin{equation}
          \label{eq:2}
          \alpha(s) = \left( \frac{a}{c} \int\sin\theta(s) ds, \frac{a}{c} \int\cos\theta(s) ds, \frac{b}{c} s \right),
          \end{equation}
          where \( c^2 = a^2 + b^2 \), is a helix, and that \( \frac{k}{\tau} = \frac{a}{b} \).
          \begin{proof}
            We need to show that there is some \( \hat{a} \) such that \( \langle T, \hat{a} \rangle = \text{constant} \).
            Additionally, we need to show that \( \frac{k}{\tau} = \frac{a}{b} \).

            To begin, note that
            \begin{equation}
            \label{eq:45}
            T(s) = \alpha^{\prime}(s) = (\frac{a}{c}\sin\theta(s), \frac{a}{c}\cos\theta(s), \frac{b}{c}),
            \end{equation}
            from which it is clear that \( \langle T, \hat{z} \rangle = \frac{b}{c} = \text{constant}\), so \(\alpha\) is a helix.

            Next, we calculate \( k \) and \( \tau \).
            \begin{align}
            \label{eq:46}
              k(s) &= \left| \alpha^{\prime\prime} \right| = \left| (\frac{a}{c} \theta^{\prime}(s) \cos\theta(s), -\frac{a}{c} \theta^{\prime}(s) \sin\theta(s), 0 ) \right| \\
              &= \sqrt{\frac{a^2}{c^2}(\theta^{\prime}(s))^2} = \frac{a}{c} \theta^{\prime}(s).
            \end{align}

            Now, since \(\tau\) is defined to satisfy \( B^{\prime} = \tau N \), we first find the binormal vector.
            \begin{align}
            \label{eq:48}
              B &= T \times N \\
                &= (\frac{a}{c}\sin\theta(s), \frac{a}{c}\cos\theta(s), \frac{b}{c}) \times (\cos\theta(s), -\sin\theta(s), 0) \\
              & = (\frac{b}{c}\sin\theta(s), \frac{b}{c}\cos\theta(s), - \frac{a}{c}).
            \end{align}

            From \cref{eq:48}, we see that
            \begin{equation}
            \label{eq:49}
            B^{\prime} = (\frac{b}{c}\theta^{\prime}(s) \cos\theta(s), -\frac{b}{c}\theta^{\prime}(s) \sin \theta(s), 0),
            \end{equation}
            therefore \( \tau = \frac{b}{c} \theta^{\prime}(s) \).

            Our claim now follows since
            \begin{equation}
            \label{eq:50}
            \frac{k}{\tau} = \frac{\frac{a}{c} \theta^{\prime}(s)}{\frac{b}{c}\theta^{\prime}(s)} = \frac{a}{b}.
            \end{equation}
          \end{proof}
  \end{enumerate}
\end{problem}
\end{document}
